% P16: density-free Markov construction and the general mediation bridge.
\section{Measurable causal graphs and their global Markov property}
\label{sec:p16-measurable-causal-graphs}

Lauritzen formulates the equivalence of directed Markov properties using
recursive probability kernels, separately from the special case of
densities \cite[Theorems~1.28, 2.44 and 2.49; Propositions~2.46--2.48,
pp.~19--20, 65--68]{Lauritzen2019GraphicalLectures2020}.
This section supplies the complete standard Borel construction and closes
the general-space graphical premise in the mediation treatment. Its graphs
have finitely many vertices; their coordinate spaces need not be finite,
countable, or dominated by a common product measure.

\subsection{Conditional independence as a kernel identity}

\begin{hbdefinition}{def:p16-measurable-ci}{Conditional independence of random elements}
Let \(X,Y,Z\) be random elements in nonempty standard Borel spaces on
\((\Omega,\mathcal F,P)\). Write \(X\perp Y\mid Z\) when
\[
 E[u(X)v(Y)\mid Z]=E[u(X)\mid Z]E[v(Y)\mid Z]
\]
almost surely for every bounded measurable real \(u,v\).
The sigma-algebra in \(E[\cdot\mid Z]\) is \(\sigma(Z)\).
For vertex sets, this notation applies to their indexed tuples; the empty
tuple is a constant. Regular conditional laws mean the common probability
kernels of \cref{sec:p15-regular-conditional-laws}.
\end{hbdefinition}

\begin{hbtheorem}{thm:p16-measurable-ci-calculus}{proved=annex:p16.measurable-ci-calculus}{Common product kernels and the semigraphoid operations}
Conditional independence is equivalent to both of the following:
\begin{enumerate}
\item \(E[u(X)\mid Y,Z]=E[u(X)\mid Z]\) for all bounded measurable \(u\).
\item A regular conditional law of \((X,Y)\) given \(Z=z\) equals
the product of its two conditional marginal laws, outside one null set
of \(z\), simultaneously as probability measures.
\end{enumerate}
For any additional standard Borel random element \(W\), it satisfies
symmetry, decomposition, weak union and contraction:
\[
\begin{aligned}
 X\perp(Y,W)\mid Z&\Longrightarrow X\perp Y\mid Z,\\
 X\perp(Y,W)\mid Z&\Longrightarrow X\perp Y\mid(Z,W),\\
 X\perp Y\mid Z,\quad X\perp W\mid(Y,Z)
     &\Longrightarrow X\perp(Y,W)\mid Z.
\end{aligned}
\]
Measurable functions may be applied to either independent tuple, also
using the conditioning tuple as an argument. No positivity, intersection
axiom or assumption of a density is required.
\end{hbtheorem}

\subsection{A finite graph of probability kernels}

\begin{hbdefinition}{def:p16-measurable-dag-kernels}{An everywhere specified kernel model}
Let \(G=(N,E)\) be a finite DAG and let \(S_v\) be a nonempty standard
Borel space for every \(v\in N\). Assign a probability kernel
\(K_v:S_{\operatorname{pa}(v)}\leadsto S_v\), specified at every parent
argument. Products of coordinate spaces carry their product sigma-algebras.
The empty coordinate product is a singleton; the empty graph carries its
point-mass law.
Starting in a topological order, successively combine the current joint
measure with \(K_v(x_{\operatorname{pa}(v)},dx_v)\) using
\cref{thm:p15-kernel-joint-measure}. Denote the resulting law by \(P_K\).
The notation
\[
 P_K(dx_N)=\mathop{\prod}_{v\in N}^{\mathrm{top}}
                         K_v(x_{\operatorname{pa}(v)},dx_v)
\]
means this iterated kernel construction, not multiplication of densities.
A local kernel on a parent stratum absent under \(P_K\) is still part of
this model. An observed probability law alone does not specify that choice.
\end{hbdefinition}

\begin{hbtheorem}{thm:p16-measurable-kernel-construction}{proved=annex:p16.measurable-kernel-construction}{Existence, order independence and ancestral marginals}
The iterated construction gives a unique probability measure and does not
depend on the topological order. Nonnegative measurable functions admit
the corresponding iterated integral; signed functions do so when absolutely
integrable. If \(D\subseteq N\) is ancestral, its marginal is the kernel
model on \(G[D]\) with exactly the retained kernels. In particular every
topological prefix has that marginal, and
\(K_v(X_{\operatorname{pa}(v)},\cdot)\) is a conditional law of \(X_v\)
given all its predecessors. All assertions hold for singular laws.
\end{hbtheorem}

\begin{hbtheorem}{thm:p16-measurable-kernel-randomization}{proved=annex:p16.measurable-kernel-randomization}{Jointly measurable realization of an entire kernel}
For any probability kernel \(K:T\leadsto S\), with \(S\) nonempty
standard Borel and \(T\) measurable, there is a jointly measurable
\(F:T\times(0,1)\to S\) such that \(F(t,U)\) has law \(K(t,\cdot)\)
for every \(t\), when \(U\) is uniform on \((0,1)\).
Consequently every finite standard Borel kernel model has a representation
\[
 X_v=F_v(X_{\operatorname{pa}(v)},U_v)
\]
with mutually independent uniform \(U_v\). This is an existence result;
the representation, and its joint coupling of different intervention
worlds, need not be unique.
\end{hbtheorem}

\begin{hbdefinition}{def:p16-measurable-markov-properties}{Ordered, local and global Markov properties}
For a topological order let \(\operatorname{pr}(v)\) be the predecessors
of \(v\), and let \(\operatorname{nd}(v)\) exclude both \(v\) and its
strict descendants. The ordered property is
\[
 X_v\perp X_{\operatorname{pr}(v)\setminus\operatorname{pa}(v)}
               \mid X_{\operatorname{pa}(v)}.
\]
The local property replaces predecessors by nondescendants. The global
property asserts \(X_A\perp X_B\mid X_C\) whenever the disjoint sets
\(A,B,C\) satisfy d-separation in \(G\). Paths, colliders and ancestral
moral separation have exactly their earlier graph-theoretic definitions.
\end{hbdefinition}

\begin{hbtheorem}{thm:p16-measurable-global-markov}{proved=annex:p16.measurable-global-markov}{Global, local, ordered and recursive-kernel equivalence}
For a probability law on the finite product of the declared standard
Borel spaces, the global, local and ordered Markov properties are
equivalent to existence of an adapted recursive kernel representation.
In particular every law \(P_K\) is globally Markov. The proof handles
a conditioned terminal vertex explicitly using weak union and contraction;
it never conditions by dividing a joint density.
\end{hbtheorem}

\subsection{Structural interventions and response backgrounds}

\begin{hbdefinition}{def:p16-measurable-scm}{A measurable structural model and interventions}
For the same finite DAG, let \(U_v\) be mutually independent random
elements in nonempty standard Borel spaces with specified laws, and let
\(f_v(x_{\operatorname{pa}(v)},u_v)\) be jointly measurable total functions.
An intervention on a fixed set \(I\) replaces its equations by constants
\(t\in S_I\). All other functions and the joint background law remain
fixed. A history-based stochastic intervention may instead use a specified
probability kernel on earlier coordinates, with independent new randomness.
The augmented graph includes exactly the permitted earlier inputs of
those replacement kernels, and is required to remain acyclic.
\end{hbdefinition}

\begin{hbtheorem}{thm:p16-measurable-interventions}{proved=annex:p16.measurable-interventions}{Measurable solutions and intervened kernel laws}
Every such model and intervention has a unique solution. For constant
interventions that solution is jointly measurable in \((t,U)\), and its
law is a probability kernel in \(t\). Its law is the recursive kernel
construction with the acted-on kernels replaced by \(\delta_{t_v}\).
History-based replacement kernels also admit the independent-noise
construction and obey the global Markov property of the resulting graph.
Consistency and composition at attained values hold pointwise. Structural
interventions are defined even when the prescribed event had zero
observed probability.
\end{hbtheorem}

\begin{hblemma}{thm:p16-measurable-response-component}{proved=annex:p16.measurable-response-component}{A separated response component gives a uniform response family}
Consider a family of measurable structural models with the same finite
DAG \(H\), same independent disturbances, and a standard Borel parameter
\(t\) in a nonempty parameter space entering their jointly measurable functions.
Let \(Y\) be a response vertex and \(B,C\) disjoint vertex sets excluding
\(Y\). Suppose \(B,C\) are pointwise unchanged as \(t\) varies and \(Y\) is d-separated
from \(B\) given \(C\) in \(H\). Expose each disturbance as a private
root. In the ancestral moral graph for \(Y,B,C\), remove \(C\) and
take the component containing \(Y\). Let \(R\) be the tuple of private
disturbances in that component. Then there is one jointly measurable
function \(F\) with
\[
 Y^t=F(t,C,R)\quad\hbox{pointwise for every }t,
 \qquad R\perp B\mid C.
\]
Thus one conditional product law of \((R,B)\) given \(C\) licenses
integration of this entire family at a measurable random parameter
depending on \((B,C)\). Separate exceptional sets for each parameter
are not substituted along a diagonal.
\end{hblemma}

\subsection{The general-space mediation graph corollary}

\begin{hbdefinition}{def:p16-measurable-mediation-data}{Graph, observations and common support measures}
Use the measurable structural model above. Let \(A,Y\) be distinct
observed vertices, with \(Y\) real-valued, and \(M\) a nonempty set of
other observed intermediate vertices, each a descendant of \(A\) and
an ancestor of \(Y\). Let observed covariate sets \(W_0,W_1,W_2,W_3\)
exclude \(A,M,Y\). Assume exactly the four separations in
\cref{thm:p16-mediation-four-set-criterion}, with
\(W_0,W_1,W_3\) nondescendants of \(A\) and \(W_2\) nondescendants
of \(M\). For the formulas, remove from \(W_1,W_3\) vertices already
in \(W_0\), and from \(W_2\) vertices already in \(W_0\cup W_1\).
This leaves each of the four conditioning sets unchanged. Put
\(C=W_0\), \(S=(W_0,W_1)\), and \(T=(A,W_0,W_1,W_2)\).
Fix exposure values \(a,b\). Define the jointly evaluated variables
\(M_b,Y_{a,m},Y_{a,M_b}\) using the intervention construction.

Let \(\nu\) be the law of \(C\), and choose regular observed kernels
\[
\begin{aligned}
 \gamma_1(c,dw_1)&=\mathcal L(W_1\mid C=c),\\
 \gamma_3(c,dw_3)&=\mathcal L(W_3\mid C=c),\\
 \gamma_2^a(c,w_1,dw_2)&=\mathcal L(W_2\mid A=a,C=c,W_1=w_1),\\
 \rho_a(c,w_1,w_2,dm)&=\mathcal L(M\mid A=a,C=c,W_1=w_1,W_2=w_2),\\
 Q^{\mathrm{obs}}_a(c,w_1,w_2,m,dy)
   &=\mathcal L(Y\mid A=a,C=c,W_1=w_1,W_2=w_2,M=m).
\end{aligned}
\]
The covariate reductions make the input tuples have distinct coordinates.
Require positive conditional exposure atoms
\[
 P(A=a\mid C,W_1)>0\quad\nu\gamma_1\hbox{-almost everywhere},
 \qquad P(A=b\mid C,W_3)>0\quad\nu\gamma_3\hbox{-almost everywhere}.
\]
The candidate reference-mediator kernel is
\[
 \eta_b(c,dm)=\int\mathcal L(M\mid A=b,C=c,W_3=w_3)(dm)
                                                    \gamma_3(c,dw_3).
\]
Require
\[
 \eta_b(c,\cdot)\ll\rho_a(c,w_1,w_2,\cdot)
 \quad \nu\gamma_1\gamma_2^a\hbox{-almost everywhere}.        \tag{MS}
\]
Every kernel denotes a probability measure at each argument, with arbitrary
measurable defaults at null observation arguments. Condition (MS) transfers
the relevant common almost-everywhere equalities to the target measure.
\end{hbdefinition}

\begin{hbtheorem}{thm:p16-measurable-mediation-fields}{proved=annex:p16.measurable-mediation-fields}{The graph supplies the uniform premises, not merely pointwise independence}
Under the stated structural graph and its nondescendant/separation
conditions, there are standard Borel disturbance tuples \(R_0,R_1,R_2\)
and jointly measurable response functions such that:
\begin{enumerate}
\item \(Y_{a,m}=F_0(a,m,C,R_0)\) for all \(a,m\), and
      \(R_0\perp M_b\mid C\) for each fixed \(b\).
\item \(Y_{a,m}=F_1(a,m,S,R_1)\) for all \(a,m\), and
      \(R_1\perp A\mid S\).
\item \(Y_m=F_2(m,T,R_2)\) for all \(m\), and
      \(R_2\perp M\mid T\).
\end{enumerate}
Each independence is an equality of conditional joint measures outside
one conditioning-null set. The first statement gives the uniform
cross-world condition of \cref{thm:p16-mediation-kernel-formula}; the
second and third license the two different observational replacements.
Also \(M_b\perp A\mid(C,W_3)\), so the displayed observed kernel
\(\eta_b\) is a version of \(\mathcal L(M_b\mid C)\) under the
exposure support condition.
\end{hbtheorem}

\begin{hbtheorem}{thm:p16-measurable-mediation-corollary}{proved=annex:p16.measurable-mediation-corollary}{The exact general mediation integral under the source graph criterion}
Under the graph, observation and support hypotheses above, every bounded
measurable \(\varphi\) satisfies
\[
\begin{aligned}
 E[\varphi(Y_{a,M_b})]
 =\int\!\int\!\int\!\int\!\int
 &\varphi(y)Q^{\mathrm{obs}}_a(c,w_1,w_2,m,dy)\eta_b(c,dm)\\
 &\gamma_2^a(c,w_1,dw_2)\gamma_1(c,dw_1)\nu(dc).
\end{aligned}
\]
This identifies the entire nested response law. For the mean require
\(E|Y|<\infty\) and finiteness of the same displayed absolute response
integral. Applying the result to the exposure pairs needed in the
definitions identifies the corresponding natural direct and indirect
effects. Only a finite union of conditioning-null sets is needed for
the finitely many chosen pairs.

The variables and mediator may be continuous, discrete, mixed, or
deterministically related. A fixed exposure point with no conditional
atom is outside this observational corollary's stated support domain;
no arbitrary conditional version at that point is claimed identifiable.
The structural intervention and the experimental kernel formula still
exist there. This extends the earlier finite graphical corollary without
assuming a product dominating measure or pointwise-to-random-index substitution.
\end{hbtheorem}

\subsection{Singular examples and the boundary of observational support}

\begin{hbtheorem}{thm:p16-measurable-graph-examples}{proved=annex:p16.measurable-graph-examples}{A singular Markov model and a continuous-mediator calculation}
Let \(U\) be uniform on \((0,1)\), and define \(X=U,Z=X,Y=Z\)
on the chain \(X\to Z\to Y\). Its law is globally Markov and is
concentrated on the diagonal. It is not absolutely continuous with
respect to any sigma-finite product of three coordinate measures.
Nevertheless its kernels and all fixed structural interventions exist.

For a mediation example take independent \(A\), a fair bit, and
uniform \(U\in(0,1)\), with
\[
 M=\begin{cases}U,&A=0,\\\sqrt U,&A=1,\end{cases}
 \qquad Y=A+M+AM.
\]
The mediator laws have densities \(1\) and \(2m\) on \((0,1)\),
and the response is deterministic conditional on its parents. The general
corollary applies with empty covariate sets and \(a_0=0,a_1=1\), and gives
\[
 \theta_{00}=1/2,\quad\theta_{01}=2/3,
 \quad\theta_{10}=2,\quad\theta_{11}=7/3.
\]
Thus PNDE is \(3/2\), TNIE is \(1/3\), and TE is \(11/6\).
Point-mediator observational identification is a different question from
this supported integral over mediator values.

Even a strictly positive observed density need not identify a fixed
non-atomic exposure point. For independent uniform \(A,M\in(0,1)\)
and an independent Bernoulli\((1/4)\) bit \(E\), compare
\(Y=E\oplus\mathbf1_{\{A=1/2\}}\) with \(Y=E\).
The same permitted mediation graph satisfies the four separations and
the observed laws agree with positive densities. Their natural direct
effects from \(1/4\) to \(1/2\) are respectively \(1/2\) and zero.
Thus a point-density assertion cannot replace the exposure-atom or
other explicitly justified identifying premises.
\end{hbtheorem}

The standard Borel and finite-graph hypotheses identify the exact
existence and graphical scope proved here. Arbitrary measurable target
spaces, infinite graphs, cyclic equations, and unrestricted conditioning
at null exposure values are separate mathematical problems.
